\documentclass{article}
\usepackage{iclr2026_conference,times}

\usepackage{hyperref}
\usepackage{url}
\usepackage{booktabs}
\usepackage{graphicx}
\usepackage{subcaption}
\usepackage{amsmath}
\usepackage{amssymb}
\usepackage{float}

\title{Match-Quality Prediction for Ranked 5v5 Games:\\
An Honest MMR + ML Pipeline}

\author{William Yu}

\newcommand{\fix}{\marginpar{FIX}}
\newcommand{\new}{\marginpar{NEW}}

\begin{document}
\maketitle

\begin{abstract}
We build and evaluate an end-to-end prototype that estimates pre-match quality
for ranked 5v5 online games. The deployed pipeline collects 500 live
\texttt{RANKED\_SOLO\_5x5} matches from the Riot Games API, reconstructs
per-player approximate MMR from public tier and LP, and assembles an
eight-dimensional match-quality feature vector including the closed-form
TrueSkill team win probability. The vector is fed to one of three
interchangeable classifiers (Random Forest, XGBoost, Logistic Regression)
that predict a three-class label --- \emph{Close}, \emph{Competitive},
\emph{Stomp} --- derived from a worst-of-three vote over duration, gold,
and kill differentials. On a second, independently-seeded 500-match
held-out collection, all three classifiers land near $50\%$ accuracy
with weak Spearman correlation ($|\rho|\!\le\!0.07$) against outcome
proxies. The agreement across model families points to a label/feature
ceiling rather than a model-capacity ceiling. We attribute the gap
primarily to (i) label noise in the heuristic
\emph{Close/Competitive/Stomp} voting rule and (ii) lossy reconstruction of
Riot's hidden MMR from public tier$+$LP.
Code and reproducibility artifacts are available at
\url{https://github.com/Onlymdestiny/Elo-Sim}.
\end{abstract}

\section{Introduction}  

Competitive online games rely on hidden rating systems (Elo, Glicko,
TrueSkill) to pair players so that matches stay close. When matchmaking
misjudges skill the result is a \emph{stomp}: one team snowballs early,
the game ends in under 22 minutes, and both sides report a poor
experience. Matchmaker designers therefore need a tool that, given the
ten pre-match rating vectors, predicts whether a match will be close,
competitive, or lopsided.

The intended user is the matchmaking engineer. The tool must run before
the match starts, use only data the server has at lobby time, and return
a fairness score together with a coarse quality label. Post-match
telemetry leaks the outcome and is not available at lobby time, so we
restrict features to lobby-time data and keep the pipeline transparent
enough to inspect and retrain.

Our contributions are (1) a full data-to-prediction prototype with a
factory boundary between feature extraction and classifier; (2) an
authenticated Riot API collection pipeline with a dual-window token-bucket
limiter, tier-balanced seed sampling, and explicit piecewise MMR
reconstruction from public tier$+$LP; and (3) a held-out evaluation on an
independently-seeded second collection, on which the three classifiers
converge near $50\%$ accuracy. Since the ceilings we observe come from
the data pipeline, Sections~\ref{sec:retrieval} and~\ref{sec:labels}
receive the same formal treatment as the rating and classifier layers.

\section{Related Work}

Rating systems have been studied since Elo~\citep{elo1978rating} and refined
into Bayesian extensions by TrueSkill~\citep{herbrich2006trueskill}, whose
closed-form team win probability we adopt directly as one feature component.
Most production matchmakers optimise a
skill-difference loss; few publish the \emph{quality-of-match} predictor
separately. Work on LoL-specific prediction has typically used post-match
statistics (gold graph, KDA) which leak outcome information and are unusable
at lobby time~\citep{do2021predicting}. Our features are strictly pre-match.

Our three-classifier factory (bagging / boosting / linear baseline) follows
the standard practice of comparing families at different bias/variance
points to distinguish a model-capacity ceiling from a feature-signal
ceiling~\citep{caruana2006empirical}.

\section{Methodology}

The pipeline has three stages: a Riot API retrieval layer producing the
training and held-out splits, a feature-extraction and labelling layer,
and a pluggable classifier factory.

\subsection{Team skill summary and win probability}
\label{sec:rating}
Each player $i$ is represented by a Gaussian skill summary $\mathcal{N}(\mu_i,
\sigma_i^2)$, where $\mu_i$ is set from the reconstructed MMR
$\widehat{M}$ defined below in \eqref{eq:mmr} and $\sigma_i$ from the
uncertainty estimate \eqref{eq:sigma}. Team-level aggregates are
\begin{equation}
\mu_T = \sum_{i \in T} \mu_i, \qquad
\sigma_T^2 = \sum_{i \in T} \sigma_i^2, \qquad T \in \{A, B\}.
\label{eq:team-agg}
\end{equation}
We use the closed-form TrueSkill team win probability as one component of
the feature vector:
\begin{equation}
p_{\text{win}}(A) \;=\; \Phi\!\left(\frac{\mu_A - \mu_B}{c}\right),
\qquad c = \sqrt{n_A \beta^2 + n_B \beta^2 + \sigma_A^2 + \sigma_B^2},
\label{eq:pwin}
\end{equation}
where $\Phi$ is the standard-normal CDF, $n_A = n_B = 5$, and $\beta = 25/6$
is the performance-variance constant. Equation \eqref{eq:pwin} is the only
rating-theoretic quantity used at inference time; we do not run an online
TrueSkill posterior update on the real data path. A simulation sandbox
that iterates the full Bayesian update was built during development and
is described in Appendix~\ref{app:sandbox}.

\subsection{Match-quality feature extraction}
\label{sec:features}
For a given pair of teams we compute the \emph{baseline} eight-dimensional
feature vector
\begin{equation}
\boldsymbol{\phi}_{\text{base}} \;=\;
\big[\,\mu_A,\; \mu_B,\; |\mu_A - \mu_B|,\;
\mathrm{Var}\{\mu_i\}_{i\in A},\; \mathrm{Var}\{\mu_i\}_{i\in B},\;
\sigma_A,\; \sigma_B,\; p_{\text{win}}\,\big] \in \mathbb{R}^8,
\label{eq:phi-base}
\end{equation}
where the team intra-variances are taken over the five constituent players.
An \emph{extended} 14-D vector appends six real-CSV-only features
\begin{equation}
\boldsymbol{\phi}_{\text{ext}} = \big[\boldsymbol{\phi}_{\text{base}},\;
\bar{w}_A,\; \bar{w}_B,\; |\bar{w}_A - \bar{w}_B|,\;
\Delta h,\; \Delta f,\; \Delta u\big],
\label{eq:phi-ext}
\end{equation}
with $\bar{w}_T$ the team-mean solo-queue winrate and
$\Delta h, \Delta f, \Delta u$ the team differences in Riot's
\texttt{hotStreak}, \texttt{freshBlood}, and \texttt{inactive} counts.

\subsection{Real-data retrieval pipeline}
\label{sec:retrieval}
We collect live \texttt{RANKED\_SOLO\_5x5} (queue 420) matches through the
public Riot API, via three subsystems.

\paragraph{(a) Tier-balanced seed sampling.}
A naive single-seed BFS collapses to one MMR cluster (friends play friends).
Let $\mathcal{T} = \{\text{Iron}, \ldots, \text{Diamond}, \text{Master},
\text{GM}, \text{Challenger}\}$ and let $S_t \subset \mathcal{P}$ denote the
set of PUUIDs returned for tier $t$. We sample uniformly
$|S_t| \le k$ per tier (default $k = 100$) from the league-v4 division or
apex endpoints and interleave by a random permutation $\pi$:
\begin{equation}
\mathcal{S} \;=\; \pi\!\left(\bigcup_{t\in\mathcal{T}} S_t\right),
\qquad \mathbb{E}[|\mathcal{S}|] \le k|\mathcal{T}|.
\label{eq:seeds}
\end{equation}
This guarantees empirical MMR coverage across the full ladder.

\paragraph{(b) Dual-window token-bucket limiter.}
Riot development keys enforce two simultaneous caps:
$N_{\text{s}} = 20$ requests per $W_{\text{s}} = 1\,$s and
$N_{\ell} = 100$ requests per $W_{\ell} = 120\,$s. Let
$Q_{\text{s}}(t), Q_{\ell}(t)$ be deques of the last-request timestamps within
each window at time $t$. A request is admitted when
$|Q_{\text{s}}| < N_{\text{s}} - \delta_{\text{s}}$ \emph{and}
$|Q_{\ell}| < N_{\ell} - \delta_{\ell}$; otherwise the thread sleeps for
\begin{equation}
\tau(t) \;=\; \max\!\Big(
\underbrace{W_{\text{s}} - (t - \min Q_{\text{s}})}_{\text{short-window wait}},\;
\underbrace{W_{\ell} - (t - \min Q_{\ell})}_{\text{long-window wait}},\;
\epsilon_0 \Big),
\label{eq:ratelimit}
\end{equation}
where $\epsilon_0 = 10^{-3}\,$s prevents busy-spinning and
$\delta_{\text{s}} = 2, \delta_{\ell} = 5$ are safety buffers. A
ten-thread pool can burst up to $N_{\text{s}} - \delta_{\text{s}} = 18$
concurrent requests whenever the long window permits.

\paragraph{(c) MMR reconstruction from tier/LP.}
Riot exposes public tier $t$, division $r$ and LP $\ell$, not the hidden
matchmaking MMR $M^\star$. We approximate $M^\star$ with a piecewise-linear
mapping $\widehat{M}(t, r, \ell)$:
\begin{equation}
\widehat{M}(t, r, \ell) \;=\;
\begin{cases}
B(t) + R(r) + \ell, & t \in \{\text{Iron}, \ldots, \text{Diamond}\},\\
3200 + \ell, & t \in \{\text{Master}, \text{GM}, \text{Challenger}\},
\end{cases}
\label{eq:mmr}
\end{equation}
where $B(t)$ is a tier base ($B(\text{Iron})=400, \ldots, B(\text{Diamond})=2800$,
step $400$), $R(\text{I})=300, R(\text{II})=200, R(\text{III})=100, R(\text{IV})=0$,
and apex tiers share one continuous LP ladder above $B(\text{Master})=3200$.
The reconstruction error $M^\star - \widehat{M}$ is the primary source of
feature noise we identify in Section~\ref{sec:discussion}.

Per-player rating uncertainty is estimated from tier, games played $n$,
and winrate $w$ via
\begin{equation}
\widehat{\sigma}(t, n, w) \;=\;
\sigma_0(t) \cdot \min\!\left(1.5,\; \tfrac{50}{n+50}\right)
\cdot \left(1 + 2|w - 0.5|\right),
\label{eq:sigma}
\end{equation}
where $\sigma_0(t)$ decreases from $210$ (Iron) to $60$ (Challenger) as a
tier-specific base. The $50/(n+50)$ factor mirrors the $1/\sqrt{n}$
convergence rate of a Bayesian posterior; the winrate term inflates
uncertainty for players still climbing or falling.

\subsection{Match-quality label}
\label{sec:labels}
Each match receives a conservative \emph{worst-of-three} label derived from
three independent outcome signals: game duration $s$, end-game gold
differential $g$, and kill differential $\kappa$. Each signal emits a
categorical vote $v_s, v_g, v_\kappa \in \mathcal{L} =
\{\text{Close}, \text{Competitive}, \text{Stomp}\}$ via thresholded
piecewise-constant functions, e.g.
\begin{equation}
v_s(s) = \begin{cases}
\text{Stomp} & s < 1320 \\
\text{Close} & s > 2100 \\
\text{Competitive} & \text{otherwise}
\end{cases},
\quad
v_g(g) = \begin{cases}
\text{Stomp} & g > 10{,}000 \\
\text{Close} & g < 4{,}000 \\
\text{Competitive} & \text{otherwise}
\end{cases},
\label{eq:votes}
\end{equation}
and analogously for $v_\kappa$ with thresholds $\{15, 7\}$. Let
$\text{sev}: \mathcal{L} \to \{0,1,2\}$ with
$\text{sev}(\text{Close})=0$, $\text{sev}(\text{Competitive})=1$,
$\text{sev}(\text{Stomp})=2$. The final label aggregates by
worst-severity:
\begin{equation}
y \;=\; \arg\max_{v \in \{v_s, v_g, v_\kappa\}}\, \text{sev}(v).
\label{eq:label}
\end{equation}
A match must therefore look close on all three axes before it is labelled
\emph{Close}. On the 500-match real set the resulting distribution is 58\%
\emph{Stomp}, 38\% \emph{Competitive}, 4\% \emph{Close}, which imposes a
substantial class-imbalance constraint on downstream accuracy.

\subsection{Classifier factory and fairness score}
\label{sec:classifier}
A \texttt{build\_model(name)} factory returns one of three estimators with
a uniform \texttt{fit}/\texttt{predict\_proba} interface: Random Forest
(class-weighted, 100 trees), XGBoost (multinomial softprob,
sample-weighted via $\text{compute\_sample\_weight}$ for imbalance) and
Logistic Regression inside a \texttt{StandardScaler} pipeline
(class-weighted, L2-regularised). All three implement the same contract
$f_\theta: \mathbb{R}^d \to \Delta^{|\mathcal{L}|-1}$ mapping a feature
vector to a class-probability simplex.

Given $\mathbf{p} = f_\theta(\boldsymbol{\phi}) = (p_C, p_M, p_S)$ where
$C, M, S$ index Close / Competitive / Stomp, we derive a continuous
fairness score on $[0, 100]$ as
\begin{equation}
\mathrm{Fairness}(\mathbf{p}) \;=\;
\big\lfloor 100\,(\,p_C + \tfrac{1}{2} p_M\,) + 0.5 \big\rfloor.
\label{eq:fairness}
\end{equation}
The deployed endpoint returns $\mathrm{Fairness}(\mathbf{p})$; the
discrete label $\arg\max_k p_k$ is kept as a diagnostic output.

\section{Experiments}
\label{sec:exp}

We evaluate the deployed pipeline
(\eqref{eq:mmr}, \eqref{eq:sigma}, \eqref{eq:phi-base}, \eqref{eq:fairness};
no online rating update) by \emph{train-on-real, test-on-real-held-out}:
train on one 500-match collection and test on a second,
independently-seeded 500-match collection gathered on a different day.
Metrics are label accuracy against the voted ground truth, the
majority-class baseline (always-\emph{Stomp}: 58\%), and Spearman $\rho$
between the continuous fairness score and each outcome proxy (Spearman
rather than Pearson because fairness and duration/gold/kill live on
different scales).

\begin{table}[t]
\caption{Held-out real-data evaluation. All three classifiers converge to
$\approx$50\% accuracy --- below the 58\% majority baseline --- with weak
rank correlation. Adding winrate and Riot status-flag differences (extended)
does not help.}
\label{tab:results}
\begin{center}
\small
\begin{tabular}{lccccc}
\toprule
 & \textbf{Features} & \textbf{Accuracy} & $\rho_{\text{dur}}$ & $\rho_{\text{gold}}$ & $\rho_{\text{kill}}$ \\
\midrule
Majority baseline & --       & 0.58 & --     & --     & --     \\
Random guess      & --       & 0.33 & --     & --     & --     \\
\midrule
Random Forest     & baseline & 0.50 & $+$0.06 & $-$0.05 & $-$0.04 \\
XGBoost           & baseline & 0.49 & $+$0.07 & $-$0.06 & $-$0.05 \\
Logistic Reg.     & baseline & 0.51 & $+$0.06 & $-$0.05 & $-$0.04 \\
\midrule
Random Forest     & extended & 0.48 & $+$0.05 & $-$0.05 & $-$0.04 \\
XGBoost           & extended & 0.47 & $+$0.06 & $-$0.06 & $-$0.05 \\
Logistic Reg.     & extended & 0.50 & $+$0.06 & $-$0.05 & $-$0.04 \\
\bottomrule
\end{tabular}
\end{center}
\end{table}

\begin{figure}[t]
\centering
\begin{subfigure}[b]{0.48\textwidth}
  \includegraphics[width=\linewidth]{real_xgb_baseline_validation_confusion.png}
  \caption{Confusion matrix (XGB, baseline).}
  \label{fig:confusion}
\end{subfigure}\hfill
\begin{subfigure}[b]{0.48\textwidth}
  \includegraphics[width=\linewidth]{real_xgb_baseline_validation_fairness_hist.png}
  \caption{Predicted fairness-score distribution by true label.}
  \label{fig:fairhist}
\end{subfigure}

\vspace{0.4em}

\begin{subfigure}[b]{\textwidth}
  \includegraphics[width=\linewidth]{real_xgb_baseline_validation_scatter.png}
  \caption{Fairness score vs.\ the three outcome proxies (duration, gold diff, kill diff).}
  \label{fig:scatter}
\end{subfigure}
\caption{Held-out real-data evaluation (XGBoost, baseline features). The
predicted fairness distribution in (b) barely separates by true label,
and (c) shows no visible trend --- consistent with the weak Spearman
$\rho$ in Table~\ref{tab:results}.}
\label{fig:pipeline}
\end{figure}

\subsection{Three-family convergence implies a feature ceiling}
Random Forest (high-variance bagging), XGBoost (sequential boosting) and
Logistic Regression (linear / low-variance) sit at different points on the
bias-variance curve. If model capacity were the limiting factor, at least
one family would separate from the others; instead the three land within
2 percentage points on every feature set. This is the feature-ceiling
signature of \citet{caruana2006empirical}: model tuning cannot recover
signal that is not in $\phi$.

\subsection{Extended features do not help}
Adding per-team winrate and Riot's \texttt{hotStreak} /
\texttt{freshBlood} / \texttt{inactive} flags (features the server knows
at lobby time but the 8-D vector ignores) does not lift accuracy
(Table~\ref{tab:results}, lower block). Winrate correlates strongly with
tier, which MMR already encodes, so the new columns are nearly collinear;
at $n=500$ the status flags are sparse enough for tree splitters to fit
noise, which is why accuracy drops slightly for RF and XGB while the
regularised Logistic Regression stays flat.

\subsection{Where is the remaining signal?}
Two candidates dominate. The first is \textbf{label noise}: the heuristic
thresholds on duration, gold, and kill differential are reasonable but not
canonical, and two annotators would disagree on borderline games, which
places an information-theoretic ceiling on any classifier. The second is
\textbf{MMR reconstruction error}: tier+LP is a lossy projection of the
hidden MMR that the matchmaker actually uses, and two Diamond II 50 LP
accounts can differ by hundreds of internal MMR points that our pipeline
cannot distinguish. Closing either gap requires data we do not have
access to.

\section{Discussion}
\label{sec:discussion}

The prototype is usable as designed: a call to
\texttt{predict\_fairness(team\_a\_mmrs, team\_b\_mmrs, model)} returns a
calibrated score and label in milliseconds. The discriminative power of
that score is modest on real data, which matters for how it is deployed:
treating it as a high-accuracy oracle would mis-route players, whereas
using it as a soft fairness regulariser alongside an existing
skill-difference objective would not.

\paragraph{Data is the bottleneck, not the model.}
Every quantity that caps accuracy is an artifact of the data pipeline:
the reconstruction error $M^\star - \widehat{M}$ in \eqref{eq:mmr}, the
label noise from the fixed thresholds in \eqref{eq:votes}, and the class
imbalance induced by \eqref{eq:label}. None of these move with more
training data or a different classifier. The pipeline is reproducible:
the collector \eqref{eq:ratelimit}, seed sampler \eqref{eq:seeds},
reconstruction constants \eqref{eq:mmr}, and voting thresholds
\eqref{eq:votes} are module-level constants, swappable independently of
the model layer.

\section{Limitations and Future Work}
\label{sec:limitations}

\paragraph{Label definition is author-chosen, not canonical.}
The Close / Competitive / Stomp thresholds on duration, gold, and kill
differential in \eqref{eq:votes} roughly track surrender timings and
gold-graph inflection points discussed in the community, but they are not
produced by a studio or validated in a rater study. Borderline matches
would be labelled differently by two reasonable annotators, and the
worst-of-three rule in \eqref{eq:label} amplifies that disagreement: one
stretched signal flips the whole label. Re-calibrating the thresholds,
replacing the vote with a soft ordinal label, or substituting
studio-canonical quality labels would each raise the information-theoretic
ceiling.

\paragraph{The ``MMR'' we use is a remap of visible rank, not Riot's MMR.}
The feature vector consumes $\widehat{M}(t, r, \ell)$ from \eqref{eq:mmr},
a piecewise-linear function of publicly visible tier/division/LP. In
Riot's system, visible rank is downstream of a hidden MMR~$M^\star$: the
matchmaker pairs players by $M^\star$ and converts $M^\star$-based
outcomes into LP movements via gain/loss rules. Two accounts at the same
visible rank can therefore carry different $M^\star$ (e.g.\ one climbing
with a positive-MMR margin and one decaying), and the reconstruction
error $M^\star - \widehat{M}$ is correlated with exactly the information
that drives the outcome. Closing this ceiling requires access to
$M^\star$ or a stronger reconstruction that folds in signals such as
recent LP-per-win/loss, games-since-promotion, and MMR-check flags.

\paragraph{Other limitations and future directions.}
The held-out set is 1000 NA1 matches over two collection days, which does
not capture regional matchmaker-tuning or patch-to-patch meta shifts. The
three-class formulation also discards the natural ordering
Close~$\prec$~Competitive~$\prec$~Stomp. Beyond the label and MMR fixes
above, promising extensions include ordinal regression or a cumulative-link
model that respects the ordering, calibration diagnostics and conformal
prediction so \eqref{eq:fairness} ships with a coverage guarantee,
multi-region and multi-patch held-out sets, and active collection
targeting the under-represented Close and Stomp tails.

\section{Conclusion}

We built an end-to-end match-quality prototype covering Riot API
collection, MMR reconstruction, a formal feature and labelling layer, and
a three-model classifier factory producing the fairness score
\eqref{eq:fairness}. A held-out evaluation on an independent 500-match
collection lands near $50\%$ accuracy, and the three model families
agree closely enough to locate the bottleneck at label noise and MMR
reconstruction rather than at model choice. For a matchmaker team, that
points to canonical quality labels and internal-MMR access as higher-yield
investments than further classifier tuning;
Section~\ref{sec:limitations} enumerates concrete next steps.

\subsubsection*{LLM Acknowledgments}
\textbf{This report is partially generated with the assistance of a large
language model (Claude, Anthropic).} Specifically, the LLM was used to
(i) draft initial prose for the Abstract, Introduction, Related Work, and
Discussion sections, which the author then edited for accuracy and scope;
(ii) typeset equations and tables from author-supplied specifications;
(iii) provide code review on the rate-limiter and feature-extraction
pipeline; and (iv) help organise the LaTeX source. All experimental
design, methodology decisions, data collection, numerical results, plots,
and final editorial control remain the author's responsibility, and every
LLM-drafted passage was reviewed and revised before inclusion. The LLM
did not have access to the private Riot API key or to the collected match
data at any point.

\bibliographystyle{iclr2026_conference}
\begin{thebibliography}{9}
\bibitem[Caruana \& Niculescu-Mizil(2006)]{caruana2006empirical}
Caruana, R. and Niculescu-Mizil, A.
\newblock An empirical comparison of supervised learning algorithms.
\newblock In \emph{ICML}, 2006.

\bibitem[Do et al.(2021)]{do2021predicting}
Do, T.~D. et al.
\newblock Using machine learning to predict game outcomes based on player-champion experience in League of Legends.
\newblock In \emph{FDG}, 2021.

\bibitem[Elo(1978)]{elo1978rating}
Elo, A.~E.
\newblock \emph{The Rating of Chessplayers, Past and Present}.
\newblock Arco, 1978.

\bibitem[Herbrich et al.(2006)]{herbrich2006trueskill}
Herbrich, R., Minka, T., and Graepel, T.
\newblock TrueSkill: A Bayesian skill rating system.
\newblock In \emph{NeurIPS}, 2006.

\end{thebibliography}

\appendix
\section{Appendix: Dataset and reproducibility details}

The 500-match real dataset was collected on 2026-04-19 from the NA1
platform, Americas routing. Seed sampling used
\texttt{--players-per-tier 100} across Iron II, Silver II, Gold II,
Platinum II, Emerald II, Diamond II, and the apex by-queue rosters
(Master / Grandmaster / Challenger). Collected PUUIDs are de-duplicated
against a persistent \texttt{.seen\_matches.json} cache so subsequent
runs extend rather than overlap the dataset.

All CSV schema, label-voting thresholds, MMR base/offset tables, and
classifier hyperparameters are surfaced as module constants at the top of
their respective files; no values are hardcoded inside loops. The
auto-incrementing output filename (\texttt{real\_matches.csv} $\to$
\texttt{real\_matches2.csv} $\to$ \ldots) makes multi-run collection safe
by default.

\section{Appendix: Simulation sandbox (design-phase tool, not deployed)}
\label{app:sandbox}

During early prototype development we built a simulation harness that
iterates the full Bayesian TrueSkill posterior update after each synthetic
match. Given $t = \mu_A - \mu_B$ and draw margin $\varepsilon$, winners
and losers are pushed via the standard truncated-Gaussian update
\begin{equation}
\mu_i \leftarrow \mu_i + \frac{\sigma_i^2}{c}\, v\!\left(\tfrac{t}{c},
\tfrac{\varepsilon}{c}\right),
\qquad
\sigma_i^2 \leftarrow \sigma_i^2 \cdot \left[1 - \frac{\sigma_i^2}{c^2}\,
w\!\left(\tfrac{t}{c}, \tfrac{\varepsilon}{c}\right)\right],
\label{eq:update}
\end{equation}
with factors
$v(t, \varepsilon) = \phi(t-\varepsilon)/\Phi(t-\varepsilon)$ and
$w(t, \varepsilon) = v(t,\varepsilon)\big(v(t,\varepsilon)+t-\varepsilon\big)$.
Simulated players are initialised with $\mu_i \sim
\max(400, \mathcal{N}(2800, 1300^2))$ and $\sigma_i \sim
\mathcal{U}(60, 210)$.

This sandbox was useful for sanity-checking the feature-extraction and
labelling code against a controllable rating distribution. It was
\emph{not} used to generate any number reported in
Section~\ref{sec:exp}: all held-out results come from the real Riot
dataset, and the deployed prediction path (equations \eqref{eq:mmr},
\eqref{eq:sigma}, \eqref{eq:phi-base}, \eqref{eq:fairness}) never invokes
\eqref{eq:update}. We document the sandbox here only for reproducibility
of the development process.

\end{document}
